\documentclass[a6paper,
              paper=landscape,
              12pt,
              pagesize,
              typearea,
              titlepage,
              DIV=10,
              twoside
              ]{scrbook}


\usepackage[utf8]{inputenc}
\usepackage[T1]{fontenc}
\usepackage[ngerman]{babel}
              

%Zusätzliche Beschnittzugabe zum DinA6 Format dazu
\usepackage[cam, noinfo, width=15.4truecm,height=11.1truecm, center]{crop}

%Überschriftenabstände
\usepackage{titlesec}

%Kopf- und Fußzeilen
\usepackage{fancyhdr}

%Schrift
\usepackage[sfdefault]{AlegreyaSans}

%Gitarrenakkorde
\usepackage{guitar}

%Index
\usepackage{makeidx}

%Strophennummerierung
\usepackage{enumitem}

%PDF-Dateien einbinden 
\usepackage{pdfpages}

%Layoutgeraffel
%may the latex masters be gentle and hope that they won't take this clusterf**k personal - but it works
%setting up the layout spacing 
\areaset[10mm]{130 mm}{82 mm}
\topmargin=-21.4mm 
\headheight=5.5mm
\headsep=2.5mm
\footskip=6mm
\setlength{\parindent}{0em}
\setlist[enumerate]{topsep=0pt, itemsep=5pt, partopsep=-10pt, parsep=5pt, wide, labelwidth=!, labelindent=0pt, label=\bfseries{\arabic*.}}
\setcounter{secnumdepth}{0}

\setlength\parskip{1\baselineskip plus 1\baselineskip minus 0.15\baselineskip}
\titleformat{\chapter}[display]  
{\normalfont\huge\bfseries}{\chaptertitlename\ \thechapter}{15pt}{\Huge}   
\titlespacing*{\chapter}{0pt}{-15mm}{5mm}
\titlespacing*{\section}{0mm}{0mm}{-5mm}
\fancypagestyle{plain}{\fancyhf{}%
  \renewcommand{\headrulewidth}{0pt}
}
\renewcommand{\footrulewidth}{0.4pt}% default is 0pt
\renewcommand{\headrulewidth}{0.4pt}% default is 0pt

%Befehlesetup
\newcommand{\whzl}{\raisebox{-1mm}{\hspace{0.3mm}\rule{0.3mm}{4mm}\hspace{0.3mm}\rule{0.1mm}{4mm}\hspace{0.2mm}}:}
\newcommand{\whzr}{:\raisebox{-1mm}{\hspace{0.2mm}\rule{0.1mm}{4mm}\hspace{0.3mm}\rule{0.3mm}{4mm}}}
\newcommand{\liedweiter}{ \AddToShipoutPictureFG*{%
                            \AtTextLowerLeft{%
                              \makebox[\textwidth][r]{
                                  \includegraphics[scale=0.025]{icons/wegzeichen_naechste.pdf}
                              }%
                            }%
                          }
                          \newpage 
}
\newcommand{\LiedSetup}{ 
                        \noindent
                        \section{\Titel}
                        \fancyhead[RE,LO]{\nouppercase{\Titel}}
                        \fancyhead[RO,LE]{\small{\nouppercase{\Interpret}}} 
                        \fancyfoot[LE,RO]{\small{\nouppercase{\Referenz}}} 
                        \fancyfoot[LO,RE]{\thepage}
}
\newcommand{\GeistlichSetup}{
                        \noindent
                        \section{\Titel}
                        \fancyhead[RE,LO]{\nouppercase{\Titel}}
                        \fancyhead[RO,LE]{\small{\nouppercase{\Interpret}}} 
                        \fancyfoot[LE,RO]{\small{\nouppercase{\Referenz}}} 
                        \fancyfoot[LO,RE]{\thepage}
                        \addcontentsline{lof}{section}{\Titel}
}                            

\makeindex

\begin{document}
  \pagestyle{empty}
  

    %Cover 
    \includepdf[scale=1.1]{bilder/klappentext_modern.pdf}
    %leerer Klappentext
    \cleardoublepage
    
    %Hier ein Vorwort, welches nicht im Inhaltsverzeichnis auftaucht
    \chapter*{Vorwort}
      Hi,

      ich bin das Vorwort von dem Liederbuchtemplate das du gerade geöffnet hast.
      Hier könntest du dich allen zukünftigen Lesern deines Liederbuches vorstellen und mitteilen was für großartige Abende sie damit verbringen können.
      
      Gezeichnet,

      dein Liederbuchtemplate

      \newpage

    %Das Inhaltsverzeichnis  
    \tableofcontents
    \newpage
 

    %Seitenrahmen einrichtung
    \pagenumbering{arabic}
    \fancyhf{}
    \pagestyle{fancy}
    
    % Dieses Template enhält nur ein einziges Lied. Wegen Liedrechten und solchen Dingen
    \def\Titel{Die Gedanken sind frei}
\def\Interpret{Deutsches Volkslied}
\def\Referenz{Möglicher Querverweis auf ein gebräuchliches Liederbuch deiner Wahl}

% hier entweder \LiedSetup() oder \GeistlichSetup{} wenn es im Kuratenverzeichnis auftauchen soll
\LiedSetup{}

\begin{guitarMagic}
	\begin{enumerate}

        \item Die Ge[A]danken sind frei, wer [E7]kann sie er[A]raten? \\
        Sie fliehen vorbei wie [E7]nächtliche [A]Schatten.\\
        Kein [E]Mensch kann sie [A]wissen, kein [E]Jäger er[A]schießen
        mit [D]Pulver und [A]Blei.\\
        Die Ge[E]danken sind [A]frei!

        \item Ich [A]denke, was ich will und [E]was mich be[A]glücket,\\
        doch alles in der Still’
        und [E]wie es sich [A]schicket.\\
        Mein [E]Wunsch und Be[A]gehren
        kann [E]niemand ver[A]wehren,
        es [D]bleibet da[A]bei:\\
        Die Ge[E]danken sind [A]frei!
        % So wird ein Seitenumbruch mit dem "auf der nächsten Seite gehts weiter" eingefügt
        \liedweiter{}

        \item Und [A]sperrt man mich ein
        im [E7]finsteren [A]Kerker,\\
        das alles sind rein
        ver[E7]gebliche [A]Werke,\\
        denn [E]meine Ge[A]danken
        zer[E]reißen die [A]Schranken
        und [D]Mauern ent[A]zwei.\\
        Die Ge[E]danken sind [A]frei!

        \item Drum [A]will ich auf immer
        den [E7]Sorgen ent[A]sagen\\
        und will mich auch nimmer
        mit [E7]Grillen mehr [A]plagen.\\
        Man [E]kann ja im [A]Herzen
        stets [E]lachen und [A]scherzen
        und [D]denken da[A]bei:\\
        Die Ge[E]danken sind [A]frei!

        \item Ich [A]liebe den Wein,
        mein [E7]Mädchen vor [A]allen,\\
        sie tut mir allein
        am [E7]besten ge[A]fallen.
        
        \liedweiter{}

        Ich [E]bin nicht al[A]leine
        bei [E]meinem Glas [A]Weine,
        mein [D]Mädchen da[A]bei.\\
        Die Ge[E]danken sind [A]frei!

    \end{enumerate}
\end{guitarMagic}

% Bilder können nicht in der guitarmagic Umgebung eingefügt werden. Das muss davor oder danach geschehen
\begin{picture}(0,0)
    \put(-5,-150){\hbox{\includegraphics[scale=0.25, angle=0]{bilder/Flaschenpost.png}}}
\end{picture}  
    

    \pagestyle{empty}

    %Das Kuratenverzeichnis fasst alle entsprechend markierten Lieder in einem weiteren Inhaltsverzeichnis zusammen
    \renewcommand\listfigurename{Kuratenverzeichnis}
    \addcontentsline{toc}{chapter}{\listfigurename{}}
    \listoffigures
    \printindex


    %Leerer hinterer Klappentext
    \cleardoublepage
    %Letzte Seite
    \includepdf[scale=1.1]{Bilder/klappentext_modern.pdf}
    



    \end{document}
    